\section{$K_0$}

\begin{definition}
$R$を環とし$P$を$R$加群とする.\ 任意の全射準同型$f \colon M \twoheadrightarrow N$と任意の準同型$g \colon P \to N$に対して次の図式が可換になる準同型$h \colon P \to M$が存在するとき$P$は$R$上の\textbf{射影加群}(projective module)であるという.
% https://q.uiver.app/#q=WzAsMyxbMSwwLCJQIl0sWzAsMSwiTSJdLFsxLDEsIk4iXSxbMSwyLCJmIiwyLHsic3R5bGUiOnsiaGVhZCI6eyJuYW1lIjoiZXBpIn19fV0sWzAsMiwiZyJdLFswLDEsInteXFxleGlzdHN9aCIsMix7InN0eWxlIjp7ImJvZHkiOnsibmFtZSI6ImRhc2hlZCJ9fX1dXQ==
\[\begin{tikzcd}
	& P \\
	M & N
	\arrow["{{^\exists}h}"', dashed, from=1-2, to=2-1]
	\arrow["g", from=1-2, to=2-2]
	\arrow["f"', two heads, from=2-1, to=2-2]
\end{tikzcd}\]

\end{definition}

射影加群の特徴づけを与える.
\begin{proposition}
$R$を環とし$P$を$R$加群とする.\ 次の(1)と(2)は同値.
\def\labelenumi{(\theenumi)}
\begin{enumerate}
    \item $P$は$R$上の射影加群.
    \item 任意の全射準同型$r \colon M \twoheadrightarrow P$に対してある準同型$s \colon P \to M$が存在して$r \circ s = \id_P$となる.
\end{enumerate}
\end{proposition}

\begin{proof}
(1)$\Rightarrow$(2): $P$を$R$上の射影加群とし$r \colon M \twoheadrightarrow P$を全射準同型とすると$\id_P \colon P \to P$は準同型なので次の図式を可換にする準同型$s \colon P \to M$が存在する.
% https://q.uiver.app/#q=WzAsMyxbMSwwLCJQIl0sWzAsMSwiTSJdLFsxLDEsIlAiXSxbMSwyLCJyIiwyLHsic3R5bGUiOnsiaGVhZCI6eyJuYW1lIjoiZXBpIn19fV0sWzAsMiwiXFxtYXRocm17aWR9X1AiXSxbMCwxLCJ7XlxcZXhpc3RzfXMiLDIseyJzdHlsZSI6eyJib2R5Ijp7Im5hbWUiOiJkYXNoZWQifX19XV0=
\[\begin{tikzcd}
	& P \\
	M & P
	\arrow["{{^\exists}s}"', dashed, from=1-2, to=2-1]
	\arrow["{\mathrm{id}_P}", from=1-2, to=2-2]
	\arrow["r"', two heads, from=2-1, to=2-2]
\end{tikzcd}\]
(2)$\Rightarrow$(1): 任意の全射準同型$r \colon M \twoheadrightarrow P$に対してある準同型$s \colon P \to M$が存在して$r \circ s = \id_P$となると仮定する.\ $f \colon M \twoheadrightarrow N$を全射準同型とし$g \colon P \to N$を準同型とする.\ $E = \{(m,p) \in M \oplus P \mid f(m) = g(p)\}$とすると$E$は$M \oplus P$の部分$R$加群である.\ $\pi_2 \colon E \to P$を$\pi_2(m,p) = p$で定めると任意の$p \in P$に対して$g(p) \in N$であり$f \colon M \twoheadrightarrow N$は全射なので$f(m) = g(p)$となる$m \in M$が存在し$\pi_2(m,p) = p$となる.\ よって$\pi_2 \colon E \to P$は全射である.\ 仮定より$\pi_1 \circ s = \id_P$を満たす$s \colon P \to E$が存在する.\ $p \in P$に対して$s(p) = (m,p) \in E$となる$m \in M$が存在する.\ $\pi_1 \colon E \to M$を$\pi_1(m,p) = m$で定め$h = \pi_1 \circ s$とする.\ 任意の$p \in P$に対して$(h(p), p) \in E$なので$f(h(p)) = g(p)$なので$f \circ h = g$である.
\end{proof}